\documentclass[ctcc_full_documentation.tex]{subfiles}
\begin{document}
\chapter{Master Code Map}

\section{Introduction}

This chapter maps every step of CTCC's execution: from entry (CLI or API) through argument parsing, environment setup, script list construction, sequential run of calculation scripts, and final BCR and output (CSV or JSON). A full workflow diagram in Mermaid syntax is available in \texttt{docs/ctcc\_execution\_workflow.txt}.

\section{Execution Flow Diagram}

\begin{center}
\resizebox{\textwidth}{!}{%
\begin{tikzpicture}[scale=0.85,
    node distance=0.45em and 0.7em,
    box/.style={rectangle, draw, fill=gray!8, rounded corners, minimum height=0.45cm, font=\small},
    smallbox/.style={rectangle, draw, fill=gray!12, rounded corners, font=\footnotesize},
    arrow/.style={-{Stealth[length=1.8mm]}, semithick}
]
\node[box] (entry) {Entry: CLI (ctcc.py) or API (ctcc\_processor.py)};
\node[box, below=0.5em of entry] (parse) {Parse CLI args (--capital\_only, --no\_insurance, etc.)};
\node[box, below=0.45em of parse] (load) {Load BCR config (YAML/JSON); override args if present};
\node[box, below=0.45em of load] (env) {Set env vars (CTCC\_INPUT\_MODE, etc.)};
\node[box, below=0.45em of env] (scen) {Generate or use CTCC\_SCENARIO\_ID};
\node[box, below=0.5em of scen] (cap) {Build script list: capital\_only (4) or full (with conditionals)};
\node[box, below=0.5em of cap] (loop) {For each script: subprocess in scripts/ with env};
\node[smallbox, below=0.35em of loop] (sload) {Script: load via smart\_loaders (YAML or JSON)};
\node[smallbox, below=0.28em of sload] (scalc) {Script: perform calculations};
\node[smallbox, below=0.28em of scalc] (swrite) {Script: write via smart\_output (CSV or JSON)};
\node[box, below=1em of swrite] (branch) {All scripts done; check CTCC\_OUTPUT\_MODE};
\node[smallbox, below left=0.8em and 0.4em of branch] (csv) {CSV: BCR from batch\_summary.csv; update CSV};
\node[smallbox, below right=0.8em and 0.4em of branch] (json) {JSON: aggregate json\_output\_*.json; write ctcc\_results\_\textlangle id\textrangle.json};
\node[box, below=1.6em of branch] (out) {Output: outputs/};

\draw[arrow] (entry) -- (parse);
\draw[arrow] (parse) -- (load);
\draw[arrow] (load) -- (env);
\draw[arrow] (env) -- (scen);
\draw[arrow] (scen) -- (cap);
\draw[arrow] (cap) -- (loop);
\draw[arrow] (loop) -- (sload);
\draw[arrow] (sload) -- (scalc);
\draw[arrow] (scalc) -- (swrite);
\draw[arrow] (swrite) -- (branch);
\draw[arrow] (branch) -| (csv);
\draw[arrow] (branch) -| (json);
\draw[arrow] (csv) |- (out);
\draw[arrow] (json) |- (out);
\end{tikzpicture}%
}
\end{center}

\section{Every Step (Numbered)}

\begin{enumerate}[leftmargin=*]
    \item Entry: user runs \texttt{python ctcc.py} (CLI) or POSTs to API \texttt{/api/ctcc/calculate} (server writes combined JSON to temp file and calls \texttt{ctcc.py} via subprocess).
    \item Parse command-line arguments: \texttt{--norisk}, \texttt{--simple}, \texttt{--no\_wf\_liability}, \texttt{--no\_emissions}, \texttt{--no\_linelosses}, \texttt{--capital\_only}, \texttt{--no\_oandm}, \texttt{--no\_insurance}, \texttt{--no\_delay\_costs}, \texttt{--no\_wildfire}, \texttt{--no\_outages}, \texttt{--no\_congestion}, \texttt{--no\_curtailment}, \texttt{-j}, \texttt{-o}, \texttt{--id}, etc.
    \item Load primary BCR config from YAML or JSON (based on \texttt{CTCC\_INPUT\_MODE}); if present, merge with CLI flags (config overrides).
    \item Build \texttt{BCRConfig} from args and set environment variables (\texttt{CTCC\_NO\_CONGESTION}, \texttt{CTCC\_NO\_WF\_LIABILITY}, etc.).
    \item Generate \texttt{CTCC\_SCENARIO\_ID} if not set (e.g. \texttt{YYYYMMDD\_HHMMSS\_microseconds}), or use existing (e.g. from \texttt{--id} or API).
    \item If JSON input mode: set \texttt{CTCC\_JSON\_DATA\_FILE} (and \texttt{os.environ}) to the combined data file path.
    \item Build script list: if \texttt{--capital\_only}, list is \texttt{weighted\_miles.py}, \texttt{build\_costs.py}, \texttt{row\_costs.py}, \texttt{environmental\_mitigation.py}. Otherwise, start with \texttt{weighted\_miles.py}, \texttt{build\_costs.py}, \texttt{row\_costs.py}, \texttt{environmental\_mitigation.py}, \texttt{revenue.py}; then append \texttt{insurance\_costs.py} unless \texttt{--no\_insurance}, \texttt{delay\_costs.py} unless \texttt{--no\_delay\_costs}, \texttt{wildfire\_costs.py} unless \texttt{--no\_wildfire}, \texttt{outage\_costs.py} unless \texttt{--no\_outages}; then \texttt{congestion\_curtailment\_reduction.py}, \texttt{energy\_losses.py}; then \texttt{oandm.py} unless \texttt{--no\_oandm}; then \texttt{emissions.py} unless \texttt{--no\_emissions}, \texttt{line\_loss\_costs.py} unless \texttt{--no\_linelosses}.
    \item For each script in order: run \texttt{sys.executable script\_name} as subprocess with \texttt{cwd=scripts}, \texttt{env=os.environ.copy()}; capture stdout/stderr; track success/failure.
    \item Each script: (a) loads data via \texttt{smart\_loaders} (dispatches to YAML or JSON loaders per \texttt{CTCC\_INPUT\_MODE}); (b) performs its calculations; (c) writes results via \texttt{smart\_output} (CSV or JSON per \texttt{CTCC\_OUTPUT\_MODE}).
    \item After all scripts: if no successful runs, exit with error. Otherwise, if \texttt{CTCC\_OUTPUT\_MODE=csv}: run \texttt{calculate\_and\_display\_bcr} on \texttt{batch\_summary.csv}, update \texttt{batch\_summary.csv} with BCR metrics via \texttt{CSVOutputManager}. If \texttt{CTCC\_OUTPUT\_MODE=json}: \texttt{aggregate\_json\_outputs} loads and merges \texttt{json\_output\_\textlangle id\textrangle\_*.json}, \texttt{build\_bcr\_data\_from\_json} extracts flat dict, BCR computed via \texttt{calculate\_benefits}/\texttt{calculate\_costs}/\texttt{calculate\_bcr\_metrics}, \texttt{write\_final\_json\_output} writes \texttt{ctcc\_results\_\textlangle id\textrangle.json}, then intermediate JSON files are removed.
    \item Exit: code 0 if all scripts succeeded; 1 if any failed.
\end{enumerate}

\section{Script Inclusion Table}

Table~\ref{tab:code-scripts} shows when each script is included in the run list and its purpose in one line.

\begin{table}[htbp]
\centering
\small
\resizebox{\textwidth}{!}{%
\begin{tabular}{@{}lll@{}}
\toprule
\textbf{Script} & \textbf{When included} & \textbf{Purpose} \\
\midrule
weighted\_miles.py & Always & Terrain-adjusted miles; prerequisite for build\_costs, insurance \\
build\_costs.py & Always & Construction costs (conductor, structure, converter; AFUDC/PV) \\
row\_costs.py & Always & Right-of-way costs \\
environmental\_mitigation.py & Always & Environmental mitigation costs \\
revenue.py & Always (full run) & Revenue requirement \\
insurance\_costs.py & If not \texttt{--no\_insurance} & Insurance costs \\
delay\_costs.py & If not \texttt{--no\_delay\_costs} & Construction/delay costs \\
wildfire\_costs.py & If not \texttt{--no\_wildfire} & Wildfire risk costs \\
outage\_costs.py & If not \texttt{--no\_outages} & Outage risk costs \\
congestion\_curtailment\_reduction.py & Always (full run) & Congestion and curtailment benefits \\
energy\_losses.py & Always (full run) & Loss data; prerequisite for line\_loss\_costs \\
oandm.py & If not \texttt{--no\_oandm} & O\&M costs \\
emissions.py & If not \texttt{--no\_emissions} & Emissions costs \\
line\_loss\_costs.py & If not \texttt{--no\_linelosses} & Line loss costs/benefits \\
\bottomrule
\end{tabular}%
}
\caption{Scripts: inclusion condition and one-line purpose.}
\label{tab:code-scripts}
\end{table}

With \texttt{--capital\_only}, only the first four scripts (weighted\_miles, build\_costs, row\_costs, environmental\_mitigation) are run.

\end{document}
